% from ega4-2-fr.tex
\oldpage[IV-2]{69}

\S~7, n\textsuperscript{o}~5, déf.~1, lorsque $A$ est de rang \emph{fini} sur
$k$ (\emph{loc. cit.}, cor. de la prop.~7), mais non en général, une $k$-algèbre
pouvant avoir un radical $\neq0$ même si elle est intègre.

\begin{corollary}[4.6.3]
\label{IV.2.4.6.3-fr}
Soit $X$ un $k$-préschéma intègre~; pour que $X$ soit géométriquement réduit
(resp. géométriquement intègre) sur $k$, il faut et il suffit que son corps de
fonctions rationnelles $\mathrm R(X)$ soit une extension séparable (resp.
séparable et primaire) de $k$.
\end{corollary}

\begin{proof}
Cela résulte aussitôt de (4.5.9) et (4.6.1).
\end{proof}

\begin{corollary}[4.6.4]
\label{IV.2.4.6.4-fr}
Soit $X$ un $k$-préschéma réduit. Alors, pour toute extension séparable $k'$ de
$k$, $X'=X\otimes_k k'$ est réduit.
\end{corollary}

\begin{proof}
Il suffit d'appliquer l'équivalence de a) et e) dans (4.6.1), en remplaçant
$X$ par $\operatorname{Spec}(k')$ et $S$ par $X$.
\end{proof}

\begin{proposition}[4.6.5]
\label{IV.2.4.6.5-fr}
\begin{enumerate}
\item[(i)] Soient $X$ un $k$-préschéma, $K$ une extension de $k$. Pour que $X$
soit géométriquement réduit (resp. géométriquement intègre) sur $k$, il faut et
il suffit que $X\otimes_k K$ soit géométriquement réduit (resp. géométriquement
intègre) sur $K$.
\item[(ii)] Soient $X$, $Y$ deux $k$-préschémas. Si $X$ et $Y$ sont
géométriquement réduits (resp. géométriquement intègres) sur $k$, il en est de
même de $X\times_k Y$.
\end{enumerate}
\end{proposition}

\begin{proof}
L'assertion (i) est conséquence triviale des définitions et de (4.4.1). Pour
démontrer la partie de l'assertion (ii) concernant la séparabilité, on observe
que si $\Omega$ est une extension algébriquement close de $k$, on a
\[
(X\times_k Y)\otimes_k\Omega=X\times_k(Y\otimes_k\Omega)~;
\]
comme $Y\otimes_k\Omega$ est réduit, il en est de même de
$X\times_k(Y\otimes_k\Omega)$ en vertu de (4.6.1, a)). Le reste de l'assertion
(ii) résulte de ce qui précède et de (4.5.8, (ii)).
\end{proof}

\begin{proposition}[4.6.6]
\label{IV.2.4.6.6-fr}
Soit $X$ un $k$-préschéma de type fini. Il existe une extension radicielle finie
$k'$ de $k$ telle que $(X_{(k')})_{\mathrm{red}}$ soit géométriquement réduit
sur $k'$.
\end{proposition}

\begin{proof}
Comme $X$ peut être recouvert par un nombre fini d'ouverts affines $U_i$, on
peut se borner au cas où $X$ est affine~: en effet, si pour chaque $i$, $k_i$
est une extension radicielle finie de $k$ telle que
$(U_i\otimes_k k_i)_{\mathrm{red}}$ soit géométriquement réduit sur $k_i$, on
peut supposer que les $k_i$ sont contenus dans une même extension radicielle
finie $k'$ de $k$~; $(U_i\otimes_k k')_{\mathrm{red}}$ s'identifie à
\[
(U_i\otimes_k k_i)_{\mathrm{red}}\otimes_{k_i}k'
\]
en vertu de (I, 5.1.8), donc est géométriquement réduit sur $k'$ par hypothèse,
et il en est donc de même de $(X\otimes_k k')_{\mathrm{red}}$. Supposons donc
que $X=\operatorname{Spec}(A)$, où $A$ est une $k$-algèbre de type fini, et
soit $\Omega$ une clôture algébrique de $k$. Désignons par $\mathfrak N'$ le
nilradical de $A\otimes_k\Omega$~; comme $A\otimes_k\Omega$ est un anneau
noethérien, $\mathfrak N'$ est un idéal engendré par un nombre fini d'éléments
de la forme
\[
y_i=\sum_j x_{ij}\otimes\xi_{ij},
\qquad x_{ij}\in A,\quad \xi_{ij}\in\Omega.
\]
Soient $K$ une sous-extension finie de $\Omega$ contenant les $\xi_{ij}$, et
$\mathfrak N$ l'idéal de $A\otimes_k K$ engendré par les $y_i$
($A\otimes_k K$ étant identifié à un sous-anneau de $A\otimes_k\Omega$)~; il
est clair que les $y_i$ sont nilpotents dans $A\otimes_k K$~; d'autre part,
comme $\mathfrak N\otimes_K\Omega=\mathfrak N'$ (et par suite
$\mathfrak N'\cap(A\otimes_k K)=\mathfrak N$), $\mathfrak N$ contient le
nilradical de $A\otimes_k K$, donc il est égal à ce nilradical. Comme
\[
(A\otimes_k\Omega)/\mathfrak N'
=\bigl((A\otimes_k K)/\mathfrak N\bigr)\otimes_K\Omega,
\]
on voit que $(X\otimes_k K)_{\mathrm{red}}\otimes_K\Omega$ est réduit, et par
suite (4.6.1), $(X_{(K)})_{\mathrm{red}}$ est géométriquement réduit sur $K$.
En remplaçant $K$ par l'extension quasi-galoisienne sur $k$ qu'il engendre dans
$\Omega$, on voit par (I, 5.1.8) que l'on peut supposer en outre $K$
quasi-galoisienne sur $k$, donc extension séparable d'une extension radicielle
finie $k'$ de $k$ (Bourbaki, \emph{Alg.}, chap.~V, \S~10,
n\textsuperscript{o}~9, prop.~14). En vertu

